% Part: applied-modal-logics
% Chapter: epistemic-logic
% Section: properties-accessibility

\documentclass[../../../include/open-logic-section]{subfiles}

\begin{document}

\olfileid{aml}{el}{acc}

\olsection{علاقات الإتاحة والمبادئ المعرفية}

في ضوء ما نعرفه عن مطابقة الأطر في المنطقيات الموجهية العادية، قد
نرغب في معرفة صورة ال!!{formula}s المميزة عند إعطائها تأويلات معرفية.
وقد قلنا إن المنطقيات المعرفية تُفسَّر عادة في S5. فلننظر إذن في كيفية
تمثيل مبادئ معرفية مختلفة، وفي مطابقتها لشروط مختلفة على الأطر.

تذكّر من المنطق الموجهي العادي أن ال!!{formula}s الموجهية المختلفة
كانت تميز خصائص مختلفة لعلاقات الإتاحة. ويختار الجدول الآتي بعضًا منها
يقابل مبادئ معرفية بعينها.

\begin{table}[t]
    \begin{tabular}{| p{.48\textwidth} || p{.48\textwidth} |}
      \hline
      {\emph{إذا كانت $R$ \dots}} & {\emph{فإن \dots صادقة في~$\mModel{M}$:}} \\
      \hline \hline
        
      & $\Knows (p \lif q) \lif (\Knows p \lif \Knows q)$ 
      \hfill \newline{(الانغلاق)} \\
      \hline
      \emph{انعكاسية}: $\forall w Rww$  
      & $\Knows p \lif p$ \hfill \newline{(الصدقية)} \\
      \hline
      \emph{متعدية}: \newline
      $\forall u \forall v \forall w ((Ruv \land Rvw) \lif Ruw)$ & 
      $\Knows p \lif \Knows \Knows p$ \hfill 
           \newline{(الاستبطان الإيجابي)} \\
      \hline 
      \emph{إقليدية}: \newline
      $\forall w \forall u \forall v ((Rwu \land Rwv) \lif Ruv)$ &  
      $\lnot \Knows p \lif \Knows \neg \Knows p$ \hfill 
        \newline{(الاستبطان السلبي)} \\
      \hline
    \end{tabular}
    \caption{أربعة مبادئ معرفية.}
    \ollabel{tab:four}
  \end{table} 

غالبًا ما تُعامل الصدقية، المقابلة لمسلمة~$T$، بوصفها أقل هذه المبادئ
إثارة للخلاف؛ فهي تمثل دعوى أن !!a{formula} إذا كانت معلومة فلا بد أن
تكون صادقة. أما الانغلاق والاستبطانان الإيجابي والسلبي فأشد إثارة
للخلاف بكثير.

يمثل الانغلاق، المقابل لمسلمة~$K$، فكرة أن معرفة العامل مغلقة تحت
الاستلزام. وقد تبدو هذه الفكرة معقولة في بعض الحالات. فقد أعرف، مثلًا،
أنني إذا كنت في فيكتوريا فأنا في جزيرة فانكوفر. وبصرف النظر عن
سيناريوهات الشك الغريبة، أعرف فعلًا أنني في فيكتوريا، وينبغي لهذا أن
يدل أيضًا على أنني أعرف أنني في جزيرة فانكوفر. ولذلك قد يبدو الانغلاق
المنطقي لمعرفتي بديهيًا نسبيًا في هذه الحالة. لكننا لا نستقصي دائمًا
نتائج معرفتنا، وقد يؤدي ذلك إلى نتائج أقل بداهة في حالات أخرى.

يُعرف الاستبطان الإيجابي أحيانًا بمبدأ KK، ويُصاغ أحيانًا بالقول إنني
إذا كنت أعرف شيئًا، فأنا أعرف أنني أعرفه. وهو النظير المعرفي
لمسلمة~4. وبالمثل يُصاغ الاستبطان السلبي بالقول إنني إذا كنت
\emph{لا أعرف} شيئًا، فأنا أعرف أنني لا أعرفه، وهو نظير مسلمة~5.
ويبدو أن لكليهما أمثلة مضادة عادية نسبيًا، أكون فيها غير متيقن مما إذا
كنت أعرف شيئًا أعرفه في الواقع.


\end{document}
